\documentclass[11pt]{hmcpset}
\usepackage[margin=1in]{geometry}
\usepackage{amsmath,amssymb,enumerate,graphicx,geometry}

\newcommand{\ket}[1]{|#1\rangle}
\newcommand{\bra}[1]{\langle #1|}
\newcommand{\braket}[2]{\langle #1|#2\rangle}
\newcommand{\braketop}[3]{\langle #1| #2 | #3 \rangle}
\newcommand{\normsq}[1]{\left|#1\right|^2}
\newcommand{\prob}[2]{\normsq{\braket{#1}{#2}}}
\newcommand{\avg}[1]{\langle #1 \rangle}

\newcommand{\ketvec}[2]{\ket{\!{#1}\!\mathbf{{#2}}}}
\def\bi#1{\textsf{\textit{#1}}}

\name{}
\class{PHYS116}
\assignment{Homework 10}
%\duedate{}

\begin{document}
\problemlist{Townsend: 10.3, 10.7, 10.8, 10.19, 11.1, 11.4}

%===========================================================

\begin{problem}[Townsend Problem 10.3 (5\,pts)]
	
	A negatively charged pion (a spin-0 particle) is bound to a proton forming a pionic hydrogen atom. At time $t=0$ the system is in the state
	\[
	|\psi\rangle = \frac{1}{2} |1,0,0\rangle + \frac{1}{\sqrt{2}} |2,1,1\rangle + \frac{1}{2} |2,1,0\rangle
	\]
	In an external magnetic field in the $z$ direction, the Hamiltonian is given by
	\[
	\hat H = \frac{\mathbf{\hat p}^2}{2 \mu} - \frac{e^2}{|\mathbf{\hat r}|} + \omega_0 \hat L_z
	\]
	where $\mu$ is the reduced mass of the pion-proton system.
	
	\begin{enumerate}
		\item[(a)] What is $|\psi(t)\rangle $ for the system at time $t$? What is $\langle E \rangle$ at time $t$ for this state?
		
		\item[(b)] What is $\langle L_z \rangle$ and $\langle L_x \rangle$ at time $t$ for this state?
	\end{enumerate}
	
\end{problem}
\begin{solution}
	\vfill
\end{solution}
\pagebreak


%===========================================================

\begin{problem}[Townsend Problem 10.7 (5\,pts)]
	
	An electron is in the ground state of tritium, for which the
	nucleus is the isotope of hydrogen with one proton and two neutrons. A
	nuclear reaction instantaneously changes the nucleus into $^3{\rm{He}}$,
	which consists of two protons and one neutron. Calculate the probability
	that the electron remains in the ground state of the new atom. Obtain a
	numerical answer.
	
	\emph{Note:} It is sufficient to treat the mass of the proton and neutron as equal.
	
	
\end{problem}
\begin{solution}
	\vfill
\end{solution}
\pagebreak


%===========================================================

\begin{problem}[Townsend Problem 10.8 (5\,pts)]
	
	Show that there are no allowed energies $E < -V_0$ for the
	potential well
	\[
	V = \begin{cases}
	- V_{0}  & r < a  \\
	0  & r > a
	\end{cases}
	\]
	by explicitly solving the Schr\"{o}dinger equation and attempting to satisfy
	all the appropriate boundary conditions.
	
	\emph{Note:} $\ell = 0$ corresponds to the ground state. Start by solving the the radial equation for $u(r) = r R(r)$ with $\ell = 0$ for $r < a$ and $r > a$ and then, as the problem indicates, attempting to satisfy the appropriate boundary conditions, namely at $r = 0$, $r = a$, and $r$ approaching infinity.
	
\end{problem}
\begin{solution}
	\vfill
\end{solution}
\pagebreak



%===========================================================

\begin{problem}[Townsend Problem 10.19 (5\,pts)]
	
	The Hamiltonian for two spin-$\frac{1}{2}$ particles, one with mass $m_1$ and the other with
	mass $m_2$, is given~by
	\[
	\hat {H} = \frac{\hat {\bf p}_{1}^{2}}{2m_{1}} + \frac{\hat
		{\bf p}_{2}^{2}}{2m_{2}} + V_{a} ( | \hat {\bf r} |
	) + \left( \frac{1}{4} - \frac{\hat {\bf S}_{1} \boldsymbol{\cdot} \hat{\bf S}_{2}
	}{\hbar^{2}} \right)V_{b} (|\hat{\bf r}|)
	\]
	where $\hat{\bf r} = \hat{\bf r}_1 - \hat{\bf r}_{2}$ and
	\[
	V_{a} (r) =
	\begin{cases}
	0  & r < a \\
	V_{0} & r > a
	\end{cases}\qquad
	V_{b} (r) =
	\begin{cases}
	0 & r < b \\
	V_{0} & r > b
	\end{cases}
	\]
	with $b < a$ and $V_0$ very large and positive.
	
	
	\begin{enumerate}
		\item[(a)] Determine the normalized position-space energy eigenfunction for the
		ground state. What is the spin state of the ground state? What is
		the degeneracy? \textit{Note:} Take $V_0$ to be infinite.
		
		\item[(b)] What can you say about the energy and spin state of the first-excited
		state? Does your result depend on how much larger $a$ is than $b$? Explain.
	\end{enumerate}
	
	Note: What are the eigen states of the spin-dependent term in parenthesis? Each of these spin configurations can have its own position-space wave function associated with it.
	
\end{problem}
\begin{solution}
	\vfill
\end{solution}
\pagebreak

%===========================================================

\begin{problem}[Townsend Problem 11.1 (5\,pts)]

Consider a perturbation $\hat{H}_{1} = b\hat{x}^{4}$ to the
simple harmonic oscillator Hamiltonian
\[
\hat{H}_0 = \frac{\hat{p}_x^2}{2m} + \frac{1}{2} m\omega^2 \hat{x}^2
\]
This is an example of an anharmonic oscillator, one with a nonlinear
restoring force.
\begin{enumerate}
	\item[(a)] Show that the first-order shift in the energy is given by
	\[
	E_{n}^{(1)} = {\frac{{3\hslash ^{2}b}}{{4m^{2}\omega^{2}}}} ( {1 + 2n + 2n^{2}})
	\]
	\item[(b)] Argue that no matter how small $b $ is, the perturbation expansion will break
	down for some sufficiently large $n$. What is the physical reason?
\end{enumerate}


\end{problem}
\begin{solution}
	\vfill
\end{solution}
	\pagebreak
	

%===========================================================

\begin{problem}[Townsend Problem 11.4 (5\,pts)]

Calculate the first-order shift to the energy of the ground
state and first-excited state of a particle of mass $m$ in the one-dimensional
infinite square well
\[
V( {x} ) = {{{\begin{cases}
			{0}  & {0 < x < L}  \\
			{\infty }  & {\rm elsewhere}  \\
			\end{cases} }} }
\]
of (a) the constant perturbation $\hat{H}_{1} = V_{1}$ and (b) the linearly
increasing perturbation $\hat{H}_{1} = \varepsilon E_{1}^{(0)} \hat{x}/L$,
where $E_{1}^{(0)}$ is the unperturbed energy of the ground state and $\varepsilon \ll 1$.

\end{problem}
\begin{solution}
	\vfill
\end{solution}
	\pagebreak
	

%===========================================================

\end{document}

